\chapter{Conclusiones y limitaciones}
\label{ch:conclusiones}

\section{Hallazgos principales}

\begin{enumerate}[leftmargin=*]
    \item mBERT B2 alcanza F1-macro \textbf{0{,}7610} en test ($n=460$), supera a LinearSVC ($0{,}7482$) y al resto de baselines, pero \textbf{sin alcanzar} la meta de $0{,}80$ definida en la configuración del proyecto. La meta de $0{,}80$ no se alcanza ni en el punto estimado ni en el extremo superior del intervalo bootstrap.

  \item El reentrenamiento B2 mejora robustez adversarial (64/75 en batería v1
        vs.\ 45/75 del checkpoint v1; 120/144 en batería extendida) y reduce el
        sesgo dialectal (FPR gap 0.07 vs.\ 0.86), a costa de una caída puntual
        del F1-macro en test de $0{,}7777$ a $0{,}7610$.
  \item La prueba de McNemar confirma superioridad significativa de mBERT B2
        frente a MultinomialNB y frente a XLM-RoBERTa v1 ($\alpha=0{,}05$).
        Frente a LinearSVC y Logistic Regression no se rechaza la hipótesis nula.
  \item Los cortes revelan fragilidad ante ironía y sarcasmo, y peor desempeño
        en Facebook respecto a X.
    \item La sonda dialectal muestra FPR elevado en No Tóxico con jerga local (41\,\% vs.\ 16\,\%), señal de posible sesgo léxico (fuente: \texttt{reports/metrics/dialect\_bias\_probe.csv}).

    \item La batería adversarial ($n=144$, 120/144 aciertos globales) confirma brechas en robustez retórica, atajo emoji--toxicidad (grupo G) y equidad dialectal del checkpoint oficial; ver Tabla~\ref{tab:adversarial-144-families}.

  \item El benchmark zero-shot (F1-macro $0{,}14$--$0{,}31$, checkpoint v1)
        evidencia transferencia limitada fuera del dominio salvadoreño.
\end{enumerate}

\section{Limitaciones}

\begin{itemize}[leftmargin=*]
  \item Tamaño moderado del test y de algunos slices (jerga, ironía).
  \item Calibración imperfecta (ECE $\approx 0{,}16$).
  \item Benchmark externo no re-ejecutado con pesos B2.
  \item Concordancia LIME--SHAP nula en la corrida reportada; pesos globales
        LIME no exportados.
  \item El F1-macro puntual del B2 queda por debajo de LinearSVC en el train B2
        ampliado, aunque mejora en escenarios adversariales.
\end{itemize}

\section{Trabajo futuro}

\begin{itemize}[leftmargin=*]
  \item Ampliar datos etiquetados reales, especialmente ironía/sarcasmo y Facebook.
  \item Re-ejecutar benchmark zero-shot con checkpoint B2.
  \item Técnicas de calibración para reducir el ECE sin sacrificar robustez.
  \item Re-ejecutar y estabilizar el pipeline SHAP; enriquecer explicaciones
        globales LIME.
  \item Explorar modelos más recientes multilingües o específicos de español.
\end{itemize}

En conjunto, el sistema B2 es más robusto ante ataques retóricos y jerga
salvadoreña que el v1, con trazabilidad reproducible, pero requiere más datos
reales y calibración antes de un despliegue de moderación automatizada.
